\documentclass[10pt,twocolumn]{article}
\usepackage[margin=1in]{geometry}
\usepackage{amsmath,amssymb,amsthm}
\usepackage{graphicx}
\graphicspath{{figures/}}
\usepackage{booktabs}
\usepackage{hyperref}
\usepackage{xcolor}
\usepackage{microtype}
\usepackage{natbib}

\title{\textbf{aria-trl}: Practical Continual Fine-Tuning for Language Models\\
\large A Three-Mechanism Subset of ARIA, Validated Against EWC on 8 Seeds}

\author{
  Darshan Poudel\\
  Independent Researcher\\
  \texttt{poudeldarshan44@gmail.com}
}

\date{July 2026}

\begin{document}
\maketitle

% ============================================================
\begin{abstract}
% ============================================================

We present an updated evaluation of \textbf{aria-trl}, a package that adds
three continual-learning mechanisms from the ARIA architecture
\citep{poudel2026aria} — a dual fast/slow-pathway FFN, Fisher-weighted
consolidation, and per-task adapters/heads — to Hugging Face TRL's
\texttt{SFTTrainer}. This supersedes our earlier report, which compared only
against unmitigated sequential fine-tuning. Here we add a genuine
\textbf{EWC baseline} (Fisher penalty only, no fast/slow split) and, in the
course of building it, found and fixed three implementation defects in the
original mechanism: the Fisher penalty was protecting the fast pathway it was
supposed to leave free, the plasticity gate's warmup never triggered within a
task's step budget, and fresh per-task heads were cold-started with too few
steps to learn. After fixing these, we validate the result twice: once on 3
seeds used during development, and again on \textbf{5 fresh seeds held out
from all tuning}. Across all 8 seeds combined, aria-trl matches EWC on
accuracy (0.7932 vs.\ 0.7927, a difference smaller than either method's own
seed-to-seed noise) while cutting forgetting to near zero
(Backward Transfer $-0.0063$ vs.\ EWC's $-0.0320$ and unmitigated
fine-tuning's $-0.0438$), and beats unmitigated fine-tuning outright on both
metrics. We report the result honestly where it is uneven: 3 of the 8 seeds
are weak points for aria-trl, not a clean sweep.

\end{abstract}

% ============================================================
\section{Introduction}
\label{sec:intro}
% ============================================================

Our first report on aria-trl \citep{poudel2026trlv1} showed it beating
unmitigated sequential fine-tuning on a 3-task text-classification benchmark.
That comparison was true but incomplete: an unprotected baseline is a low bar,
and it says nothing about whether aria-trl's specific mechanism —
Fisher-weighted consolidation over a fast/slow pathway split — outperforms
the standard, much simpler version of that idea, Elastic Weight Consolidation
\citep{kirkpatrick2017overcoming}. This report adds that comparison, and is
honest about what building it uncovered: the original aria-trl configuration
had three interacting defects that, once fixed, changed the result
substantially. We describe the defects, the fixes, and a two-stage validation
— development seeds, then held-out seeds never used while iterating — because
a single round of seed-fitting is exactly how spurious continual-learning
results get published.

% ============================================================
\section{What Changed Since the First Report}
\label{sec:changes}
% ============================================================

\subsection{The EWC baseline}

We implemented textbook EWC as a third method: a plain \texttt{distilgpt2}
with no fast/slow pathway, no adapters, and a single shared classification
head, trained with only a Fisher-weighted penalty term. It reuses the same
\texttt{FisherConsolidator} aria-trl already had, since that component's
Fisher estimation itself was correct — the defects were elsewhere.

\subsection{Three defects, found by building a fair comparison}

\textbf{The Fisher penalty was protecting the fast pathway.} aria-trl's
fast/slow split only works if the fast lane stays free to adapt each task;
the slow lane is what should be protected. The original implementation's
Fisher penalty covered every backbone parameter, including the fast pathway
— so the "fast, freely-adapting" lane was being pulled back toward old-task
values just like everything else. Net effect: aria-trl was behaving like
EWC plus a halved effective learning rate, which is strictly worse at
learning new tasks than EWC alone.

\textbf{The plasticity gate's warmup never activated.} The gate
regularizer that pushes the fast/slow mixing weight $\pi$ toward
specializing (0 or 1, rather than blending at 0.5) only turns on after
\texttt{warmup\_steps}. That was hardcoded to 100, but at this benchmark's
scale (200 training examples, batch size 32) one task is only $\sim$13
optimizer steps — the gate never left its initial unspecialized state.

\textbf{Fresh heads and adapters were cold-started with too few steps.}
Each new task gets its own classification head and adapter, both
randomly initialized, sharing the same $\sim$13-step budget as everything
else. Standard FT and EWC, by contrast, keep one head that accumulates
training across all tasks. The undertrained heads showed up directly in the
data: aria-trl's just-learned (diagonal) accuracies were consistently below
the two baselines' on the same task.

\subsection{The fixes}

\begin{enumerate}
\item \textbf{Rescoped Fisher protection.} \texttt{FisherConsolidator} now
  explicitly excludes the fast pathway's parameters, protecting only the
  slow pathway, attention, LayerNorm, and embeddings.
\item \textbf{Warmup matched to the real step budget.} \texttt{warmup\_steps}
  is now computed as one epoch's worth of steps for the actual dataset size,
  not a fixed constant.
\item \textbf{One throttle on the slow pathway, not two.} The slow
  pathway was previously throttled by both a reduced learning rate
  (\texttt{slow\_lr\_ratio=0.5}) and gradient dampening by $(1-\bar\pi)$.
  We set \texttt{slow\_lr\_ratio=1.0}, leaving the gate-driven dampening as
  the only throttle.
\item \textbf{Higher learning rate for fresh heads and adapters.} The
  active task's head and adapter now train in their own parameter group at
  5$\times$ the base learning rate.
\item \textbf{Warm-started heads.} Each new task's head is initialized
  from the previous task's head weights rather than randomly — reasonable
  here since all three tasks are binary sentiment-adjacent and task
  identity is known at train time.
\item \textbf{Equal training budget.} All three methods now train for 5
  epochs (previously 2), so the comparison isn't confounded by aria-trl
  needing more steps than the others got.
\end{enumerate}

A fairness point worth stating plainly: aria-trl's per-task heads mean task
identity must be known at evaluation time to route to the correct head. EWC
and Standard FT share one head and do not need this. This is a legitimate
task-incremental setup, not sleight of hand, but it is a real difference in
what each method assumes, and we state it here rather than leave it implicit.

% ============================================================
\section{Experimental Setup}
% ============================================================

Same benchmark as before: SST-2 (Movies) $\to$ Yelp Review Full
(Restaurants) $\to$ dair-ai/emotion (Social), binary labels,
\texttt{distilgpt2}, 200 train / 80 test examples per task, CPU-only.
\texttt{spc\_lambda=15.0} (aria-trl and EWC), \texttt{adapter\_dim=64},
head/adapter learning-rate multiplier 5$\times$ base (2e-4). All three
methods train 5 epochs per task.

\textbf{Two-stage seed validation.} We first ran all three methods on the
same 3 seeds used throughout development (42, 123, 7) — a pilot fixes were
tuned against, so a win here alone would not be trustworthy. We then froze
the configuration completely and ran 5 seeds chosen after the fact and never
seen during development (2024, 777, 555, 99, 1234), with no further changes
permitted regardless of outcome. We report development seeds, fresh seeds,
and the combined 8 separately, rather than only the number that looks best.

% ============================================================
\section{Results}
% ============================================================

\begin{figure}[t]
\centering
\includegraphics[width=\columnwidth]{fig1_summary_bars.png}
\caption{ACC and BWT, mean $\pm$ std across all 8 seeds. aria-trl is
statistically tied with EWC on accuracy and has the best (least negative)
BWT of the three.}
\label{fig:summary}
\end{figure}

\begin{table}[t]
\centering
\footnotesize
\begin{tabular}{lccc}
\toprule
 & Std.\ FT & EWC & aria-trl \\
\midrule
ACC (dev, n=3)      & 0.7639 & 0.7903 & 0.7861 \\
ACC (fresh, n=5)    & 0.7992 & 0.7942 & 0.7975 \\
ACC (combined, n=8) & 0.7859 & 0.7927 & \textbf{0.7932} \\
\midrule
BWT (dev, n=3)      & $-$0.0563 & $-$0.0334 & $-$0.0104 \\
BWT (fresh, n=5)    & $-$0.0362 & $-$0.0313 & $-$0.0038 \\
BWT (combined, n=8) & $-$0.0438 & $-$0.0320 & \textbf{$-$0.0063} \\
\bottomrule
\end{tabular}
\caption{ACC and BWT means, reported at each validation stage rather than
only the combined figure. aria-trl's ACC is within noise of EWC at every
stage; its BWT is the best of the three at every stage.}
\label{tab:main}
\end{table}

\textbf{Combined-8 deltas:} vs.\ Standard FT, ACC $+0.0073$, BWT $+0.0375$
(aria-trl wins both outright). vs.\ EWC, ACC $+0.0005$ (a difference far
smaller than either method's own std, 0.0080--0.0178 — a statistical tie,
not a win), BWT $+0.0257$ (aria-trl forgets substantially less).

\begin{figure}[t]
\centering
\includegraphics[width=\columnwidth]{fig2_dev_vs_fresh_validation.png}
\caption{The same pattern — aria-trl tied with EWC on ACC, ahead on BWT —
holds independently in the development seeds and the fresh, previously
unseen seeds.}
\label{fig:validation}
\end{figure}

\subsection{Where it is uneven}

Backward Transfer's standard deviation for aria-trl (0.0538 combined) is
larger than EWC's (0.0280) or Standard FT's (0.0417) is not — this is not a
clean sweep. Two seeds, 777 and 1234, are simultaneous weak points where
aria-trl underperforms \emph{both} baselines on \emph{both} metrics.
Figure~\ref{fig:curves} shows all 8 seeds individually rather than hiding
this inside the aggregate: 5 of 8 seeds are clear aria-trl wins on
retention, 3 of 8 are not.

\begin{figure}[t]
\centering
\includegraphics[width=\columnwidth]{fig3_forgetting_curves_8seed.png}
\caption{Task-0 accuracy after each step, all 8 seeds shown individually.
Seeds 777 and 1234 are where aria-trl's retention advantage does not hold.}
\label{fig:curves}
\end{figure}

\begin{figure}[t]
\centering
\includegraphics[width=\columnwidth]{fig4_accuracy_matrix.png}
\caption{Full accuracy matrix, seed 42. aria-trl's off-diagonal (older-task)
values sit closer to their diagonal values than either baseline's.}
\label{fig:matrix}
\end{figure}

% ============================================================
\section{Discussion}
% ============================================================

\textbf{The accuracy race is a genuine tie, not a near-win rounded up.} A
0.0005 combined-8 ACC delta is smaller than a single flipped prediction out
of the 240 test examples evaluated per seed-batch stage. We are reporting
this as "matches EWC," not "beats EWC," because the latter would overstate
what eight seeds actually support.

\textbf{The BWT advantage is real but not uniform.} It holds in 5 of 8 seeds
and holds in the combined aggregate at every validation stage (dev, fresh,
combined) — but seeds 777 and 1234 show aria-trl losing on both metrics
simultaneously to EWC. We did not tune against these seeds after seeing
them, and we are not tuning against them now; they are reported as the
honest shape of the result, not smoothed over.

\textbf{Why fix and re-validate, rather than just report the buggy
version's numbers?} The pre-fix configuration, tested against EWC, showed
aria-trl \emph{losing} on accuracy to both baselines while winning only on
BWT — a materially different and weaker claim than the one this report
makes. We are documenting both, in this report, because a mechanism that
only wins against a strawman baseline and a mechanism that ties a real one
while forgetting less are different results, and conflating them would be
the kind of thing that shouldn't survive scrutiny.

% ============================================================
\section{Limitations}
% ============================================================

\begin{itemize}
\item \textbf{Scale.} 200 train / 80 test examples per task, CPU-only —
  chosen for iteration speed, not for representing production fine-tuning.
\item \textbf{Task-identity assumption.} aria-trl's per-task heads require
  knowing which task is active at evaluation time; EWC and Standard FT do
  not need this. Stated explicitly here (see Section~\ref{sec:changes}).
\item \textbf{Uneven BWT.} 3 of 8 seeds are simultaneous losses on both
  metrics; the aggregate advantage is real but not a guarantee per-seed.
\item \textbf{Eight seeds.} More rigorous than three, still not enough for
  a strong statistical claim beyond "the direction is consistent across two
  independently-run seed batches."
\item \textbf{Single model, single size, binary labels.} Unchanged from the
  first report: \texttt{distilgpt2} only, binary classification only.
\end{itemize}

% ============================================================
\section{Conclusion}
% ============================================================

Against a genuine EWC baseline, validated on 5 seeds held out from all
development plus the original 3, aria-trl matches EWC's accuracy within
noise and cuts forgetting to near zero — a smaller, more precisely-stated
claim than our first report's, and a more trustworthy one, since it survived
a fair comparison and an out-of-sample check rather than only an unprotected
baseline. Code, the updated benchmark script, and full per-seed results for
both validation stages are available at
\url{https://github.com/rsd-darshan/ARIA-TRL}.

% ============================================================
\bibliographystyle{plainnat}
\begin{thebibliography}{99}

\bibitem[Poudel(2026)]{poudel2026aria}
Poudel, D.
ARIA: Adaptive Recurrent Intelligence Architecture --- Morphogenic Attention
and Slow-Pathway Consolidation for Continual Learning.
2026.
\url{https://github.com/rsd-darshan/ARIA}

\bibitem[Poudel(2026)]{poudel2026trlv1}
Poudel, D.
aria-trl: Practical Continual Fine-Tuning for Language Models (first report).
2026.
\url{https://github.com/rsd-darshan/ARIA-TRL}

\bibitem[Kirkpatrick et al.(2017)]{kirkpatrick2017overcoming}
Kirkpatrick, J., Pascanu, R., Rabinowitz, N., et al.
Overcoming catastrophic forgetting in neural networks.
\emph{Proceedings of the National Academy of Sciences}, 114(13):3521--3526, 2017.

\bibitem[Buzzega et al.(2020)]{buzzega2020dark}
Buzzega, P., Boschini, M., Porrello, A., Abati, D., and Calderara, S.
Dark experience for general continual learning: a strong, simple baseline.
\emph{Advances in Neural Information Processing Systems}, 33, 2020.

\bibitem[Socher et al.(2013)]{sst2}
Socher, R., Perelygin, A., Wu, J., et al.
Recursive deep models for semantic compositionality over a sentiment
treebank.
\emph{EMNLP}, 2013.

\bibitem[Zhang et al.(2015)]{yelp}
Zhang, X., Zhao, J., and LeCun, Y.
Character-level convolutional networks for text classification.
\emph{Advances in Neural Information Processing Systems}, 28, 2015.

\bibitem[Saravia et al.(2018)]{emotion}
Saravia, E., Liu, H.C.T., Huang, Y.H., Wu, J., and Chen, Y.S.
CARER: Contextualized affect representations for emotion recognition.
\emph{EMNLP}, 2018.

\end{thebibliography}

\end{document}
